\documentclass[11pt]{article}
\usepackage{amsmath,amssymb,amsthm,hyperref}
\newtheorem{theorem}{Theorem}
\newtheorem{lemma}[theorem]{Lemma}
\newtheorem{proposition}[theorem]{Proposition}
\newtheorem{corollary}[theorem]{Corollary}
\title{P\'olya Averaging for Boolean Reciprocal Heads}
\author{}
\date{November 23, 2026}
\begin{document}
\maketitle

\paragraph{Scope.}
This is an explicit averaging construction and an audit of its local
variations. It does not prove the requested universal head bound.
All input variables below are Boolean unless specified otherwise.

\section{An exact admissible refinement identity}

Let $b_i>0$, $\delta>0$, $B(x)=1+\sum_i b_i x_i$, and
$T=1+\sum_i b_i$. Then
\[
 \frac{T}{B(x)}=\frac{1+\delta}{B(x)+\delta}
                    +\sum_{i=1}^n\frac{b_i}{B(x)+\delta x_i}.
 \tag{1}
\]
Every denominator has strictly positive constant and variable
coefficients. Thus a single reciprocal is a positive weighted average
of $n+1$ normalized reciprocal heads. The nodes are $b/(1+\delta)$
and $b+\delta e_i$, with weights $1/T$ and $b_i/T$, respectively.
They approach $b$ as $\delta$ tends to zero.

\textbf{Proof.} At a fixed input put $s=\sum_i b_i x_i=B-1$.
Grouping the summands according to $x_i=0$ or $1$ gives
\[
 \frac{1+\delta+s}{B+\delta}
                   +\frac{T-1-s}{B}
 =1+\frac{T-B}{B}=\frac TB.
\]
The algebraic identity is kernel checked as
\texttt{HeadComplexity.polya\_reciprocal\_identity} in
\texttt{Results/PolyaReciprocalIdentity.lean}.

For non-Boolean real inputs at which all denominators are nonzero,
the right side of (1) minus its left side equals
\[
 \frac{\delta^2}{B(B+\delta)}
             \sum_i\frac{b_i x_i(x_i-1)}{B+\delta x_i}.
 \tag{2}
\]
This follows by subtracting $\sum_i b_i/B$ and simplifying each
term. It is strictly negative when every $0<x_i<1$.
Consequently (1) is genuinely a Boolean-only cancellation, not a
rational identity on an open input domain.

The averaging rule is the finite conditional-expectation calculation
for a P\'olya urn: add $\delta$ units to one of weights
$(1,b_1,\ldots,b_n)$, selected proportionally to its weight. The
relative share of color zero in any selected group containing that
color is a martingale. The standard two-color calculation is in
Durrett, \emph{Probability: Theory and Examples}, version 5,
Section 4.3.2, page 226,
\url{https://sites.math.duke.edu/~rtd/PTE/PTE5_011119.pdf}.
Identity (1) is proved directly above and does not require a
probabilistic convergence theorem.

\section{The infinitesimal operator and its limit}

For a fixed input $x$, put $K_x(b)=(1+b\cdot x)^{-1}$. Direct
differentiation, using $x_i^2=x_i$, gives
\[
 \left[\sum_i b_i\partial_{ii}
       +\sum_{i,j}b_i b_j\partial_{ij}
       +2\sum_i b_i\partial_i\right]K_x=0.
 \tag{3}
\]
Indeed the first two terms sum to $2(s+s^2)/(1+s)^3$ and the
last is $-2s/(1+s)^2$, where $s=b\cdot x$. The principal matrix
$\operatorname{diag}(b)+bb^{\mathsf T}$ is positive definite:
its quadratic form at $v\ne0$ is
$\sum_i b_i v_i^2+(b\cdot v)^2>0$.
This explains why positive averaging is possible locally. The
operator identity is a manual calculation, not a Lean declaration.
Neither harmonicity nor averaging bounds the number of atoms needed
to obtain an arbitrary strict sign table.

\section{A selector formula and first variations}

For $x_i\in\{0,1\}$,
\[
 \frac1{B+\delta x_i}
   =\frac1B+x_i\left(\frac1{B+\delta}-\frac1B\right).
 \tag{4}
\]
Thus, for affine $A_i$ and a constant $q$,
\[
 \sum_i\frac{A_i}{B+\delta x_i}
 -\frac{\sum_i A_i+q}{B}+\frac q{B+\delta}
 =\left(q+\sum_i x_i A_i\right)
                \left(\frac1{B+\delta}-\frac1B\right).
 \tag{5}
\]
The two identities are kernel checked as
\texttt{HeadComplexity.boolean\_reciprocal\_selector} and
\texttt{HeadComplexity.polya\_affine\_selector\_sum}.
Any quadratic Boolean polynomial has the first factor on the right
for suitable affine $A_i$ and $q$. Its negative sign is therefore
represented by the $n+2$ admissible denominators in (5).
This restricted-class construction is not the target theorem.

A first variation of the constant-numerator identity (1), including
an arbitrary bias, has the form
\[
 c+\frac{U(x)}{B(x)^2}+\frac{V(x)}{(B(x)+\delta)^2}
       +Q(x)\left(\frac1{(B(x)+\delta)^2}-\frac1{B(x)^2}\right),
 \tag{6}
\]
where $U,V$ are affine and $Q$ is quadratic. To see this, replace
(4) by its squared-denominator version and sum the affine variation
numerators. The linear and constant parts of $Q$ can be absorbed
into $U,V$. Multiplication by $B^2(B+\delta)^2$ gives degree at
most four, because $(B+\delta)^2-B^2$ is affine.
Thus (6) cannot strictly represent parity when $n>4$.
For completeness, a polynomial of degree less than $n$ has zero
alternating sum on the Boolean cube, since every monomial misses
some coordinate. A strict parity sign representation would give a
nonzero alternating sum of one sign, a contradiction.

The exact finite compiler nevertheless produced an eight-head
certificate for the fixed seeded six-bit diagnostic target.
This checks a concrete use of (1); it is not a statement about all
six-bit functions or the asymptotic maximum.

\section{Repeated averaging and a common polynomial denominator}

After $k$ urn steps, write $\alpha=(\alpha_0,\ldots,\alpha_n)$ for
the color counts, with $\sum_j\alpha_j=k$. The resulting affine
forms are
\[
 B_\alpha(x)=B(x)+\delta\left(\alpha_0+\sum_i\alpha_i x_i\right).
 \tag{7}
\]
Repeated conditional expectation gives a positive weighted identity
for $1/B$ with these $\binom{n+k}{k}$ forms. More explicitly, put
$w=(1,b_1,\ldots,b_n)$ and
$(a)_{r,\delta}=\prod_{j=0}^{r-1}(a+j\delta)$. Then
\[
 \frac1B=\sum_{|\alpha|=k}
 \frac{k!}{\prod_j\alpha_j!}
 \frac{\prod_j(w_j)_{\alpha_j,\delta}}{(T)_{k,\delta}}
 \frac{1+\delta\alpha_0}{B_\alpha}.
 \tag{8}
\]
A given ordered sequence of colors has the product probability in
(8) without its multinomial coefficient. There are exactly that
many orderings of its counts. This proves (8) by induction from
(1), including all real positive starting weights.

\begin{lemma}[Polynomial clearing for a finite offset set]
Let $Q(t)=\prod_{j=0}^k(t+\delta j)^2$. There is a polynomial
$R(t,z)$ of total degree at most $2k$ such that
\[
       R(t,j)=Q(t)/(t+\delta j)^2\qquad(0\le j\le k).
\]
\end{lemma}
\begin{proof}
Divide $Q(t)$ by the monic polynomial $(t+\delta z)^2$, as a
polynomial in $t$ with coefficients in $\mathbb R[z]$. The quotient
has total degree at most $2k$: each long-division subtraction preserves
the total-degree bound, starting from degree $2k+2$ and removing a
monic factor of total degree two. After substituting $z=j$, the
remainder is zero because $(t+\delta j)^2$ divides $Q(t)$.
This quotient is the required $R$.
\end{proof}

\begin{proposition}[Degree bound for complete-layer first variations]
Any score
\[
                   c+\sum_{|\alpha|=k} D_\alpha(x)/B_\alpha(x)^2,
\]
with affine $D_\alpha$, has threshold degree at most $2k+2$.
This also permits a single base term with denominator $B^2$.
\end{proposition}
\begin{proof}
At a Boolean input,
$L_\alpha(x)=\alpha_0+\sum_i\alpha_i x_i$ is an integer in
$\{0,\ldots,k\}$. Multiplication by the positive $Q(B(x))$
replaces each summand by
$D_\alpha(x)R(B(x),L_\alpha(x))$, of degree at most $2k+1$.
The bias contributes degree $2k+2$. The base term uses $j=0$.
\end{proof}

Consequently a complete-layer first-variation construction with
$\binom{n+k}{k}=O(2^n/n^2)$ cannot be universal. If it represented
parity, it would require $k\ge(n-2)/2$. For $n\ge6$, this implies
$k\ge n/3$. The standard factorial estimate then gives
\[
 \binom{n+\lceil n/3\rceil}{\lceil n/3\rceil}
 =2^{((4/3)\log_2(4/3)+(1/3)\log_2 3+o(1))n}
 =2^{(1.0817\ldots+o(1))n},
\]
which exceeds $C2^n/n^2$ for every fixed $C$. This conclusion counts
the complete layer. It does not assume the same cost for a selected
sparse subset, an overlapping collection of refinements, or arbitrary
higher-order perturbations. Those possibilities need separate
constructions or invariants; no universal existence is claimed here.

\end{document}
